\scanpage{17}
\[
(\unsure{13|12})\qquad E\simeq\dot T\times[0,1]
\]
rendant commutatif
\diagram{\[
\begin{tikzcd}
\dot T\arrow[r,hook]\arrow[dr,"{t\mapsto(t,0)}"']&E\arrow[d,"\simeq",draw=none]\\
&\dot T\times[0,1]
\end{tikzcd}
\]}
et tel que $\unsure{X'{}^*|X'_*}\cap E$ s'identifie à une partie
de $\dot T\times[0,1]$ de la forme $X'_0\times[0,1]$, où $X'_0$
est une partie fermée de $\dot T$.

Soit $B$ une partie bord stricte de l'espace
topologique $X$. Considérons un vois.
cylindrique $E$ \deleted{\unsure{absolu de}}\inserted{de} $B$ dans $X$. On
dit qu'une partie \inserted{fermée} $X'$ de $X$ est \uline{compatible}
avec $E$ --- i.e. avec l'homéom.
\[
E\simeq B\times[0,1]\qquad -
\]
si $X'\cap E$ s'identifie à $B'\times[0,1]$
(où néc$^t$. on a $B'=X'\cap B$).
\margin{plutôt : \uline{transversale ?}}
\margin{Cela implique donc que $B\cap X'$ est un
sous-espace bord de $X'$.}
On fera attention que la définition de la
``\deleted{compatibilité}\inserted{transversalité}
\deleted{dans le cas} cylindrique'' \uline{n'est pas} un cas
particulier de celle de la \deleted{\unsure{en}}\inserted{compatibilité}
conique, \uline{mais est}
\inserted{plus exigeante : correspond au cas où on \unsure{lui demandait}
$\bar T'\cap Y=0.\dot T'$}
\deleted{dans ce cas $X'\cap E\simeq B'\times[0,1]\cup
B\times\{0\}\unsure{\varepsilon}$}

On définit de $=$ la compatibilité pour
une famille \deleted{finie} \deleted{\unsure{de}} parties $(X_i)_{i\in I}$.
\margin{16}
